\documentclass[10pt]{article}
\usepackage[margin=20mm]{geometry}
\usepackage{amsmath,amssymb,mathtools,fontspec,unicode-math,xurl,hyperref}
\IfFontExistsTF{DejaVu Serif}{\setmainfont{DejaVu Serif}}{\setmainfont{Latin Modern Roman}}
\setmathfont{latinmodern-math.otf}
\setlength{\parindent}{0pt}\setlength{\parskip}{6pt}\sloppy
\title{Technical Verification of
The Generalized Wehrl Entropy Bound in Quantitative Form}
\author{AI-assisted technical review}
\date{}
\begin{document}
\maketitle
Abstract: The exact 23-page arXiv version is examined through its internal proof structure. An admissible convex entropy function gives both a coherent and a noncoherent normalized pure state generalized entropy equal to negative infinity. This identifies a finiteness qualification needed for the literal uniqueness assertion in Remark 4(e), without refuting Theorem 3’s extended-real inequality. Independent checks also identify local formula corrections in the two proofs and a distance-normalization correction in the explicit pure-state constant. Their repairs are verified without asserting a counterexample to the main stability inequalities. The evidence is bound to the precise preprint version.\par
\section*{Identity and scope}
Rupert L. Frank, Fabio Nicola and Paolo Tilli, The generalized Wehrl entropy bound in quantitative form, arXiv:2307.14089v3, 1 November 2023 [source]. The associated journal DOI is 10.4171/JEMS/1674; online publication is recorded as 8 July 2025 [publisher]. Equality of this preprint and the journal version has not been established. Page references are physical PDF pages, which agree with printed pages 1–23.\par
All 23 pages of extracted text were read. Internal checks covered Proposition 6, Proposition 8, Theorem 9, the mixed-state estimates in Lemma 10, the distribution-function proof of Theorem 3, the pure-state argument in Section 3, and the Fock-space application in Section 4. Independent verification additionally read pp. 1–2 and 6–18 of the exact 46-page arXiv:2307.09304v1 of Gómez, Guerra, Ramos and Tilli [ggrt]: all five steps of Lemma 2.1, Corollary 2.2, Lemmas 2.3–2.5 and the assembly of the function-stability estimate were checked against their use here. This is not certification of that entire reference. Remark 12’s higher-dimensional extension was read but not fully verified. Source pages 1–3, 5–6, 9–14 and 18–21, and reference pages 1–2, 11–14 and 16–17, were actually inspected as bound page images.\par
\section*{Assumptions and proof structure}
Equation reference: ineq. For a positive trace-class operator \(\rho\) of trace one, the source defines \[u_\rho(z)=\sum_jp_j|F_j(z)|^2e^{-\pi|z|^2},\qquad
 D[\rho]=\inf_w\|\rho-|\varphi_w\rangle\langle\varphi_w|\|_{\mathcal S^1},\] where the Bargmann transforms have unit norm in the Fock space. Theorem 3 assumes \(\Phi:[0,1]\to\mathbb R\) continuous, convex, non-linear, and \(\Phi(0)=0\), and asserts \[
 I_\Phi(\rho)\leq I_\Phi(\rho_{\rm coh})-c_\Phi D[\rho]^2,\qquad
 I_\Phi(\rho)=\int_{\mathbb C}\Phi(u_\rho)\,dA.\] The proof on p. 12 explicitly allows negative-infinite integrals.\par
The estimate \(D[\rho]\leq2\sqrt{1-T}\), \(T=\max u_\rho\), follows from the two-dimensional pure-state trace-norm calculation, the triangle inequality and concavity of the square root. The numerator in Proposition 8 is nonnegative by Cauchy–Schwarz. Coarea and planar isoperimetry then give \(\mu'(t)\leq-1/t\) on \((0,T)\). The single crossing of \(\mu-\mu_0\) supplies the sign required in (37)–(39). The three threshold cases in the final comparison with \(H(t_*)\) were checked, including \(t_1>T\) and a non-smooth convex \(\Phi\).\par
For the imported part of Lemma 10, set \(\delta^2=(1-T)/T\) and use the conjugated harmonic term in source (25). Cauchy–Schwarz gives \[|h|\leq\sqrt2\pi\delta r^2e^{\pi r^2/2},\quad
 |h_r|\leq2\sqrt2\pi\delta r e^{\pi r^2/2},\quad
 |h_{rr}|\leq C\delta(1+r^2)e^{\pi r^2/2}.\] The last estimate follows by retaining the same coefficient factor in the second-derivative calculation; source (31) prints a valid weaker bound without \(\delta\). For fixed \(t_0>0\) and sufficiently small \(\delta\), the defining function \(g_\theta\) in (32) satisfies \(g_r\geq\pi t_0 r e^{\pi r^2}\). Its level radius \(r_\sigma\) is unique, \(\pi r_\sigma^2\leq\log(2/t_0)\), and implicit differentiation gives \(|r_\sigma'|\leq C\delta r_\sigma/t_0\) and \(|r_\sigma''|\leq C\delta^2r_\sigma/t_0^3\). For \(A(\sigma)=\tfrac12\int r_\sigma^2\,d\theta\), the harmonic mean-value identity gives \(A'(0)=0\), while \(|A''|\leq C\delta^2A(0)/t_0^3\). This includes Step V of the reference, in addition to its Steps III–IV.\par
Writing \(q=T/t\) and using \(\delta^2\leq t_0/2\), the initial radius obeys \[\pi r_0^2=\log q+\log(1-\delta^2)-\log(1-\delta^2q)
 \leq(1+2\delta^2/t_0)\log q .\] Indeed the extra logarithm is at most \(2\delta^2(q-1)\leq2\delta^2t_0^{-1}\log q\). This verifies the stated mixed-state distribution estimate. It bounds the crossing point away from one. On a fixed interval \([a,T]\) sufficiently close to one, it then gives \(\mu_0(t)-\mu(t)\geq(1-T)/2\) when \(T\) is close to one. Integration gives \(H(t_*)\geq c(1-T)\); when \(T\) stays away from one, \(H(t_*)\geq(1-T)^2/2\) supplies the same linear conclusion. Thus the relevant reference argument for Lemma 11 transfers without an additional pure-state assumption.\par
For Section 4, the monomial expansion directly verifies \[\pi^{-1}\|F'\|^2=\sum n|a_n|^2,\qquad
 \pi\int|z|^2|F|^2\,d\mu=\sum(n+1)|a_n|^2 .\] These identities justify the integration-by-parts limit when the energy is finite. Scaling and the unitary Bargmann correspondence give the asserted Fock-space implication; no sharper constant is inferred.\par
\section*{A finiteness qualification for uniqueness}
Problem 1 (Degenerate extended-real maximizers).  Remark 4(e), p. 3, says that the generalized entropy functional is maximized only by coherent states when \(\Phi\) is non-linear. The literal stated class includes a function for which this functional is negative infinite at a noncoherent state as well.\par
\subsection*{Admissibility of the entropy function}
Equation reference: phi. Set \(a=e^{-1}\) and \[
 \Phi(0)=0,\qquad
 \Phi(t)=\begin{cases}
 -(1-\log t)^{-1},&0<t\leq a,\\
 -(1+et)/4,&a\leq t\leq1 .
 \end{cases}\] For \(0<t<a\), let \(L=1-\log t>2\). Direct differentiation gives \[\Phi'(t)=-\frac1{tL^2},\qquad
 \Phi''(t)=\frac{L-2}{t^2L^3}>0.\] At \(a\), both branches have value \(-1/2\) and derivative \(-e/4\). The limit at zero is zero. Thus \(\Phi\) is continuous, convex and non-linear on \([0,1]\), with \(\Phi(0)=0\). The linear branch is essential: the first branch is not being used outside its convexity interval.\par
\subsection*{Normalized states and divergent integrals}
For the coherent state at zero, \(u_0=e^{-\pi|z|^2}\). In radial integrals put \(s=\pi|z|^2\), so that \(dA=ds\). For \(s\geq1\), \[\Phi(e^{-s})=-\frac1{1+s},\qquad I_\Phi(\rho_0)=-\infty.\] Take the Fock function \(F_1(z)=\sqrt{\pi}\,z\). Its normalization is exact: \[\int_{\mathbb C}\pi|z|^2e^{-\pi|z|^2}\,dA
 =\int_0^\infty se^{-s}\,ds=1.\] The source’s unitary Bargmann correspondence therefore supplies an admissible rank-one density matrix \(\rho_1\). Its Husimi function is \(u_1=se^{-s}\), with maximum \(e^{-1}<1\). It is noncoherent and, by Remark 7, \[D[\rho_1]=2\sqrt{1-e^{-1}}>0.\] Since \(se^{-s}\leq e^{-1}\), the first branch of [phi] (\#phi) applies. For \(s\geq1\), \[\Phi(se^{-s})=-\frac1{1+s-\log s}\leq-\frac1{1+s}.\] The integrand is nonpositive, hence this comparison proves \(I_\Phi(\rho_1)=-\infty\). No numerical truncation is being used to infer divergence.\par
\subsection*{Counterevidence, impact and exact repair}
The source’s opening paragraph on p. 12 expressly discusses \(I_\Phi(\rho_{\rm coh})=-\infty\) and says the other integral is then negative infinite as well. This is evidence that the proof does not mistake an infinite quantity for a finite one. If the uniqueness sentence is read as restricted to finite optima, its intended conclusion is consistent; that restriction is absent from the literal statement.\par
Equation [ineq] (\#ineq) remains valid as \(-\infty\leq-\infty\). The example does not refute Theorems 1 or 3 or establish a defect in the Fock-space application. Its supported impact is local, corresponding to a supported Minor qualification. The relevant source wording has now been visually confirmed. It is not a demonstrated failure of the core stability inequalities.\par
Restrict the uniqueness assertion to \[\int_0^1\frac{\Phi(t)}t\,dt>-\infty .\] The positive part is already bounded by convexity and \(\Phi(0)=0\). With a finite coherent-state optimum, equality in [ineq] (\#ineq) forces \(D[\rho]=0\) and the closed coherent-projector family gives a coherent state. This verifies the repair to the corollary conditional on Theorem 3; it is not an independent proof of that theorem.\par
\section*{Local formula corrections with complete repairs}
Problem 2 (Conjugation, ratio and negative-part notation).  Three local formulas require consistent notation. In source (22), p. 9, the coefficient \(a_0^{(j)}\) should be conjugated, as it already is in (20) and (25). On p. 18, the variable following (48) must be \(s=e^{-|E|}/\tau\), and the integrand in (49) must be \((u_0-\tau)_-\), as in (47)–(48).\par
The conjugation correction is fixed by the expansion of \(|F_j|^2\). Multiplying \(F_j\) by \(i\) preserves the density matrix but reverses the unconjugated product in (22); the conjugated expression is phase invariant. Formula (25) already supplies the correct expression used in the subsequent estimates. This local correction does not invalidate the estimates verified above.\par
The sign error in (49) has a counterexample under the actual hypotheses. Use the normalized \(F_1\) above, \(u_1=se^{-s}\), and \(\tau=e^{-3}\). For \(E=\{u_1>\tau\}\), the radial annulus \(e^{-2}\leq s\leq4\) is contained in \(E\): the values at both ends exceed \(e^{-3}\) and the function is unimodal. Thus \[|E|\geq4-e^{-2}>3+\log2=|F_\tau|+\log2,
 \qquad e^{-|E|}<\tau/2 .\] The strict inequality follows from \(e^{-2}<1/4\) and \(\log2<3/4\). The centered disk \(E^*\) therefore strictly contains \(F_\tau=\{u_0>\tau\}\). The printed integral of \((u_0-\tau)_+\) over \(E^*\setminus F_\tau\) is zero, whereas the printed lower bound is strictly positive.\par
With the negative part, \(s=e^{-|E|}/\tau\in(0,1)\) and \(\Omega=E^*\setminus F_\tau\), the exact integral is \[\int_{\Omega}(u_0-\tau)_-\,dA
 =e^{-|E|}-\tau+\tau|E|-\tau|F_\tau|
 =\tau(s-\log s-1)
 \geq\frac{\tau}{2}(1-s)^2 .\] The last inequality follows by differentiating \(s-\log s-1-(1-s)^2/2\), whose derivative is \(-(1-s)^2/s\leq0\) and whose value at one is zero. For \(s<1/2\) the remainder is at least \(\tau/8\). Since \(D[f]^2\leq2\) and the coherent hinge integral is at most one, the continuation on p. 19 follows with the claimed universal form. These are Minor local formula errors with verified repairs, not a counterexample to Theorem 13.\par
\section*{Distance normalization in the explicit constant}
Problem 3 (A missing conversion in the deduction on p. 19).  Theorem 13 measures the unit-phase \(L^2\) distance \(D[f]\). Theorem 3, equation (4), measures the trace distance \(D[\rho]\). The displayed explicit coefficient on p. 19 identifies the same constant across these two distances without the conversion.\par
For a normalized pure state \(\rho=|f\rangle\langle f|\), the identities on pp. 5–6 give \[D[f]^2=2(1-\sqrt T),\qquad
 D[\rho]^2=4(1-T)=2(1+\sqrt T)D[f]^2 .\] Let \(c\) be the actual constant in Theorem 13, and put \[A_\Phi=\int_0^1
 \tau\bigl(1-\tau-\tau\log(1/\tau)\bigr)\,\Phi''(d\tau).\] The superposition argument establishes an entropy gap at least \(cA_\Phi D[f]^2\). It therefore guarantees the coefficient \(cA_\Phi/4\) in front of \(D[\rho]^2\), rather than directly giving \(cA_\Phi\) as printed for (4). For \(\Phi(t)=t\log t\), \(A_\Phi=1/4\), so this route guarantees \(c/16\) in (2); the stated \(c/4\) is the coefficient for \(D[f]^2\).\par
This is a Minor constant-accounting defect in the explicit deduction. No counterexample to the larger numerical constant is claimed: an independently stronger estimate might exist. The existence assertions of Theorems 1 and 3 are unchanged. One can either include the factor \(1/4\) or explicitly rename the universal constant after reducing it. Theorem 14 can still use the constant from Theorem 1, because \(D[f]^2\leq D[\rho]^2\); its deduction requires an implication, not an equality of those distances.\par
\section*{Conclusion}
The independent checks support a finite-optimum qualification for Remark 4(e), local formula repairs, and an explicit distance-constant conversion. Each has a complete local repair and a supported Minor impact. The p. 12 extended-real discussion, the corrected conjugation already present in (25), and the separate first proof of Theorem 3 limit broader interpretations. No counterexample to the main stability inequalities has been established.\par
No erratum was found in the bounded exact-title search and consulted publisher/arXiv records. This is not an exhaustive priority search. Optional supplements and code were not located in the acquisition record; this is not a mathematical defect. Author code was not run.\par
Dimension | State\par
Issue judgment | Three supported local findings; see Problems 1–3\par
Supported impact | Minor; the main inequalities are not refuted\par
Repair check | Finite-optimum restriction; consistent local formulas; factor \(1/4\) in the trace-distance conversion\par
Scientific decision | Pending; no human review supplied\par
Author adoption | Not requested; unknown\par
Evidence record | Bound source images inspected; report-page checks in provenance\par
9 R. L. Frank, F. Nicola and P. Tilli, The generalized Wehrl entropy bound in quantitative form, arXiv:2307.14089v3 (1 November 2023). https://arxiv.org/abs/2307.14089v3. EMS Press, record for the same title, DOI 10.4171/JEMS/1674, online publication 8 July 2025. https://ems.press/journals/jems/articles/14298930. Consulted 27–28 October 2026; full journal PDF not acquired. J. Gómez, A. Guerra, J. P. G. Ramos and P. Tilli, Stability of the Faber–Krahn inequality for the Short-time Fourier Transform, arXiv:2307.09304v1 (18 July 2023), 46 pages. https://arxiv.org/abs/2307.09304v1. The exact v1 PDF was acquired and the focused portions specified above were read; no equivalence with another version is assumed.\par
\section*{Supporting checks}
> 1b03df278ca7493aed42ec019a613976b7da8ec070cf0460917fe23d5a431158\par
The initial reviewer-written symbolic program and its retained output were read in full. It checks derivatives, joining values and slopes, normalization, and the coherent-state improper integral. It was not rerun during independent acceptance. The present independent reasoning supplies the analytic divergence comparison, the admissible counterexample to the positive-part formula, the repaired negative-part bound, the distance conversion, and the focused proof checks above. No author code was executed and no auxiliary scientific agent was used.\par
This independent revision preserves the prior TeX, PDF and metadata. The accompanying independent build record binds the exact revised TeX, source, official Tectonic 0.17.0 executable, retained dependency payloads, actual commands, exit codes, final PDF, text extraction and page images. The page-inspection record distinguishes actual native-image observation from rendering and text parsing. Old false observations and failed image-return attempts, including `EMFILE` and incomplete transports, remain preserved; historical display times are not invented.\par
Only the assigned frozen copy was used. The original scientific staging, immutable inputs, canonical reports and controller queue were not edited. Independent scientific verification is recorded separately from formal acceptance; the scientific decision remains pending and human reviews remain empty. UI facts belong to the controller. Runtime self-description is recorded separately and is not independent backend telemetry.\par
\end{document}
